% from ega2-1-fr.tex
\oldpage[II]{148}
\subsection{Courbes algébriques et corps de fonctions de dimension 1}
\label{subsection:II.7.4-fr}

Le but de ce numéro est de montrer comment se formule en langage des schémas
la notion classique de courbe algébrique (telle qu'elle est exposée par
exemple dans le livre de C.~Chevalley [23]). Dans tout le numéro, \emph{on
désigne par $k$ un corps, tous les schémas considérés sont des $k$-schémas de
type fini, et tous les morphismes des $k$-morphismes}.

\begin{proposition}[7.4.1]
\label{II.7.4.1-fr}
Soit $X$ un préschéma de type fini sur $k$ (donc noethérien)~; soient $x_i$
$(1\leq i\leq n)$ les points génériques des composantes irréductibles $X_i$ de
$X$, et soit $K_i=k(x_i)$ $(1\leq i\leq n)$. Les conditions suivantes sont
équivalentes~:
\begin{enumerate}
\item[\rm a)] Chacun des $K_i$ est une extension de $k$ dont le degré de
transcendance est égal à 1.
\item[\rm b)] Pour tout point fermé $x$ de $X$, l'anneau local
$\mathcal{O}_x$ est de dimension 1 (\hyperref[II.7.1.5-fr]{7.1.5}).
\item[\rm c)] Les parties fermées irréductibles de $X$ distinctes des $X_i$
sont les points fermés de $X$.
\end{enumerate}
\end{proposition}

Comme $X$ est quasi-compact, toute partie fermée irréductible $F$ de $X$
contient un point fermé (\hyperref[0.2.1.3-fr]{0, 2.1.3}). En vertu de
(\hyperref[I.2.4.2-fr]{I, 2.4.2}), il y a correspondance biunivoque entre les
idéaux premiers de $\mathcal{O}_x$ et les parties fermées irréductibles de $X$
contenant $x$ (\hyperref[I.1.1.14-fr]{I, 1.1.14})~; l'équivalence de b) et c)
en découle aussitôt. D'autre part, si $\mathfrak{p}_\alpha$
$(1\leq\alpha\leq r)$ sont les idéaux premiers minimaux de l'anneau local
noethérien $\mathcal{O}_x$, les anneaux locaux
$\mathcal{O}_x/\mathfrak{p}_\alpha$ sont intègres, et ont pour corps des
fractions ceux des $K_i$ tels que $x\in X_i$. En outre, on sait ([1],
p.~4-06, th.~2) que la dimension d'une $k$-algèbre locale intègre de type fini
est égale au degré de transcendance sur $k$ de son corps des fractions. Enfin,
la dimension de $\mathcal{O}_x$ est borne supérieure des dimensions des
$\mathcal{O}_x/\mathfrak{p}_\alpha$~; or, la condition a) implique que ces
dimensions sont égales à 1, donc a) implique b)~; inversement, si
$\mathcal{O}_x$ est de dimension 1, aucun des $\mathfrak{p}_\alpha$ ne peut
être égal à l'idéal maximal de $\mathcal{O}_x$, sans quoi $\mathcal{O}_x$
serait de dimension 0~; donc chacun des
$\mathcal{O}_x/\mathfrak{p}_\alpha$ est de dimension 1, ce qui montre que b)
entraîne a).

On notera que sous les conditions de (\hyperref[II.7.4.1-fr]{7.4.1}),
l'ensemble $X$ est \emph{vide ou infini}, comme il résulte aussitôt de
(\hyperref[I.6.4.4-fr]{I, 6.4.4}).

\begin{definition}[7.4.2]
\label{II.7.4.2-fr}
On appelle courbe algébrique sur $k$ un schéma algébrique non vide sur $k$
vérifiant les conditions de (\hyperref[II.7.4.1-fr]{7.4.1}).
\end{definition}

Dans le langage de la dimension, qui sera introduit au chapitre IV, cela
s'exprimera en disant qu'une courbe algébrique sur $k$ est un $k$-schéma
algébrique non vide dont \emph{toutes les composantes irréductibles sont de
dimension 1}.

On notera que si $X$ est une courbe algébrique sur $k$, les sous-préschémas
fermés réduits $X_i$ $(1\leq i\leq n)$ de $X$ ayant pour espaces sous-jacents
les composantes irréductibles de $X$ sont des courbes algébriques sur $k$.

\begin{corollary}[7.4.3]
\label{II.7.4.3-fr}
Soit $X$ une courbe algébrique irréductible. Le seul point non fermé de $X$ est
son point générique. Les parties fermées de $X$ distinctes de $X$ sont les
ensembles finis de points fermés~; ce sont aussi les seules parties de $X$ qui
ne soient pas partout denses.
\end{corollary}

Si un point $x\in X$ n'est pas fermé, son adhérence dans $X$ est une partie
fermée irréductible de $X$, donc nécessairement $X$ tout entier en vertu de
(\hyperref[II.7.4.1-fr]{7.4.1}), et par suite $x$ est le point générique de
$X$. Une partie fermée $F$ de $X$, distincte de $X$, ne peut contenir

